\chapter{Data Description}
\label{ch:data_description}

This chapter describes the two publicly available datasets used in this thesis, their characteristics, and the preprocessing steps applied to prepare them for classification. Both datasets are sourced from Kaggle and represent symptom-based disease prediction tasks at different scales.

\section{Dataset~1: Disease Symptom Prediction Dataset}
\label{sec:dataset1}

The first dataset, the Disease Symptom Prediction Dataset \citep{Kaggle_Disease2020}, contains approximately 5,000 patient records. Each record maps a combination of symptoms to one of 41 disease categories. The dataset includes 132 unique symptoms represented as columns, where the presence of a symptom is indicated by a value of 1 and the absence by 0.

Table~\ref{tab:dataset1_overview} summarizes the key characteristics of this dataset.

\begin{table}[H]
\centering
\caption{Overview of Dataset~1: Disease Symptom Prediction Dataset}
\label{tab:dataset1_overview}
\begin{tabular}{@{}ll@{}}
\toprule
\textbf{Property}         & \textbf{Value} \\
\midrule
Source                     & Kaggle (itachi9604) \\
Number of records          & $\sim$5,000 \\
Number of features         & 132 (binary symptoms) \\
Number of disease classes  & 41 \\
Feature type               & Binary (0/1) \\
Target variable            & Disease name (categorical) \\
Missing values             & Present in raw data \\
\bottomrule
\end{tabular}
\end{table}

The dataset was originally compiled from multiple medical sources and provides supplementary files including disease descriptions, symptom severity ratings, and suggested precautions. I only used the symptom-disease mapping, since my primary objective is to evaluate classifier performance on the core prediction task.

The class distribution in Dataset~1 is moderately imbalanced. Most disease categories have between 100 and 150 records, but some rare diseases have fewer than 50 examples. I addressed this through class weighting (Section~\ref{sec:class_imbalance}).

\section{Dataset~2: Diseases and Symptoms Dataset}
\label{sec:dataset2}

The second dataset, the Diseases and Symptoms Dataset \citep{Kaggle_Diseases2023}, is substantially larger, containing approximately 246,000 patient records. It covers a broader range of conditions than Dataset~1.

Table~\ref{tab:dataset2_overview} summarizes its key characteristics.

\begin{table}[H]
\centering
\caption{Overview of Dataset~2: Diseases and Symptoms Dataset}
\label{tab:dataset2_overview}
\begin{tabular}{@{}ll@{}}
\toprule
\textbf{Property}         & \textbf{Value} \\
\midrule
Source                     & Kaggle (dhivyeshrk) \\
Number of records          & $\sim$246,000 \\
Number of features         & After preprocessing: 132 \\
Number of disease classes  & 41 \\
Feature type               & Binary (0/1) \\
Target variable            & Disease name (categorical) \\
Missing values             & Minimal \\
\bottomrule
\end{tabular}
\end{table}

The larger size of Dataset~2 lets me see how classifiers scale with more data. It also reduces the variance of performance estimates from cross-validation, which makes the comparison more reliable.

\section{Dataset Comparison}
\label{sec:dataset_comparison}

The two datasets differ in several important ways beyond size. Dataset~1 is relatively clean but small, which makes overfitting a concern, especially for complex models with many parameters. Dataset~2 provides more data but might be noisier due to how it was compiled.

I chose to use two datasets of different scales because it lets me compare classifiers not just in terms of absolute performance but also in terms of how they scale with data availability. This matters for practical deployment --- different healthcare contexts will have very different amounts of training data available.

\section{Preprocessing Pipeline}
\label{sec:preprocessing}

I applied the same preprocessing pipeline to both datasets to ensure fair comparison across classifiers. The steps were:

\textbf{Step 1: Data Loading and Inspection.} I loaded each dataset from CSV and checked for structural issues, column types, missing values, and duplicate records.

\textbf{Step 2: Missing Value Handling.} Dataset~1 had missing values in symptom columns, which I filled with 0 (symptom absent) to match the binary encoding. Dataset~2 had minimal missing values, which I handled the same way.

\textbf{Step 3: Feature Encoding.} Symptoms needed to be encoded as binary features. Dataset~1 was already in this format. For Dataset~2, I had to convert text-based symptom entries to the same binary vector representation.

\textbf{Step 4: Label Encoding.} I encoded disease names as integer labels using scikit-learn's \texttt{LabelEncoder} and kept the mapping for interpreting results later.

\textbf{Step 5: Train-Test Split.} I split each dataset into 80\% training and 20\% testing using stratified sampling to preserve the class distribution.

\textbf{Step 6: Feature Scaling.} For Logistic Regression and SVM, I standardized features using \texttt{StandardScaler} to have zero mean and unit variance. Tree-based classifiers (Random Forest, XGBoost) and Naive Bayes don't need scaling, so I trained them on the original binary features.

\textbf{Step 7: Cross-Validation Setup.} I used stratified 5-fold cross-validation for hyperparameter tuning and evaluation. Stratified splits matter for imbalanced datasets because they keep the same class proportions in each fold.

The complete preprocessing pipeline, including software versions (Table~\ref{tab:software}) and hyperparameter search spaces (Table~\ref{tab:hyperparameters}), is documented in Appendix~\ref{app:supplementary}.
